\documentclass[11pt]{article}
\RequirePackage[T1]{fontenc}
\RequirePackage[utf8]{inputenc}
\usepackage{subfiles} 


\date{}
\usepackage{color}
\usepackage{bigstrut}
\usepackage{multirow}
\usepackage{array}
\usepackage{tabularx}
\usepackage[left=2cm,top=3cm,right=2cm,bottom=2cm]{geometry}
\usepackage{setspace}
\setstretch{1.2}
\usepackage{graphicx}
\date{}
\newcommand*\ruleline[1]{\par\noindent\raisebox{.8ex}{\makebox[\linewidth]{\hrulefill\hspace{1ex}\raisebox{-.8ex}{#1}\hspace{1ex}\hrulefill}}}
\def\vhrulefill#1{\leavevmode\leaders\hrule\@height#1\hfill \kern\z@}

\newcommand{\rpm}{\sbox0{$1$}\sbox2{$\scriptstyle\pm$}
	\raise\dimexpr(\ht0-\ht2)/2\relax\box2 }
		
 \usepackage{amsthm}
 \newtheorem{lemma}{Lemma}
 \newtheorem{proposition}{Proposition}
 \theoremstyle{remark}
 \newtheorem*{remark}{Remark}		
\usepackage{mathrsfs}
\usepackage{amsmath}
\usepackage{amssymb}
\usepackage{mathtools}
\usepackage{amsmath}
\usepackage{gensymb}
%% The amssymb package provides various useful mathematical symbols
\usepackage{amssymb}
\usepackage{tikz,pgfplots}
\usepackage{dblfloatfix}

\usepackage{setspace}

\renewcommand{\baselinestretch}{1} 

%\usepackage{multibbl}
\usepackage{xcolor}
\usepackage{url}
\definecolor{newcolor}{rgb}{.8,.349,.1}

\usepackage{amsmath}
\usepackage{amsfonts}
\usepackage{amssymb}
\usepackage{multirow}
\usepackage{bigstrut}
\usepackage{graphicx}
\usepackage{subfig}
\usepackage{caption}
\usepackage{url}
\usepackage{amsfonts}
\usepackage{amsmath}
\usepackage{subeqnarray}
\newtheorem{theorem}{Theorem}
\newcommand{\sign}{\mathrm{sign}}
\newcommand{\sat}{\mathrm{sat}}
\usepackage{commath}
\newcommand*\xor{\mathbin{\oplus}}
\usepackage{xr}
\newcommand{\beginsupplement}{%
	\setcounter{table}{0}
	\renewcommand{\thetable}{S\arabic{table}}%
	\setcounter{figure}{0}
	\renewcommand{\thefigure}{S\arabic{figure}}%
}
\usepackage{pgfplots}
\usepackage{balance}
\usepackage{booktabs}
\setlength{\aboverulesep}{-0.2pt}
\setlength{\belowrulesep}{-0.2pt}

%\usepackage[capitalise, noabbrev]{cleveref}
%% Cross-reference 
% \usepackage{xr-hyper}
%\usepackage{cleveref}


%% Custom packages 
\usepackage{xcolor}
\usepackage{color}
\definecolor{mygray}{RGB}{193,203,215}
\usepackage{booktabs}
\usepackage{multirow}
\usepackage{colortbl}
\usepackage{graphicx}
%\usepackage{subcaption}
\usepackage{caption}
\DeclareCaptionFont{blue}{\color{blue}}
\usepackage{booktabs}
\usepackage{array}
\usepackage{dblfloatfix}
\usepackage{balance}
\usepackage{enumitem}
\usepackage{url}
%\setlist[enumerate]{itemsep=0mm}
\usepackage{dingbat}
\newcolumntype{C}[1]{>{\centering\arraybackslash}m{#1}}

\usepackage[framemethod=tikz]{mdframed}
%\definecolor{mycolor}{rgb}{0.9, 0.0, 0.0}
%\newmdenv[innerlinewidth=0.5pt, roundcorner=4pt,linecolor=mycolor,innerleftmargin=6pt,
%innerrightmargin=6pt,innertopmargin=6pt,innerbottommargin=6pt]{rolebox}
%\definecolor{mycolor1}{rgb}{0.0, 0.0, 0.9}
%\newmdenv[innerlinewidth=0.5pt, roundcorner=4pt,linecolor=mycolor1,innerleftmargin=6pt,
%innerrightmargin=6pt,innertopmargin=6pt,innerbottommargin=6pt]{commentbox}

\usepackage{bigstrut}
\usepackage{multirow}
\usepackage{amsmath}
\usepackage{subfigure}

\makeatletter
\long\def\@IEEEtitleabstractindextextbox#1{\parbox{0.922\textwidth}{#1}}
\makeatother


%% Pack the article in compact form
% \usepackage[compact]{titlesec}
% \titlespacing{\section}{1pt}{*0}{*0}
% \titlespacing{\subsection}{1pt}{*0}{*0}
% \titlespacing{\subsubsection}{1pt}{*0}{*0}

\usepackage[rawfloats=true]{floatrow} 
\restylefloat{figure} 
\restylefloat{table} 

\setlength{\textfloatsep}{1\baselineskip plus 0.2\baselineskip minus 0.5\baselineskip}
 %% Pack references
 

\usepackage{cite}
\usepackage{amsmath,amssymb,amsfonts}
\usepackage{algorithmic}
% \usepackage{caption,subcaption,graphicx}
\usepackage{caption,graphicx}
\usepackage{textcomp}
\usepackage{xcolor}
\usepackage{mathrsfs}
\usepackage{pgfplots}

%%%%%%%%%%%%%%
\usepackage{amssymb}% http://ctan.org/pkg/amssymb
\usepackage{pifont}% http://ctan.org/pkg/pifont
\newcommand{\cmark}{\ding{51}}%
\newcommand{\xmark}{\ding{55}}%




%%%%%%%%%%%%%%%%%
\usepackage[table,xcdraw]{xcolor}
\usepackage{url}
\usepackage{textcomp}
\usepackage{xcolor}
\usepackage{mathrsfs}
%\usepackage{subfigure}
\usepackage{booktabs} 
\usepackage{tikz}
\usepackage{dblfloatfix}
\usepackage{bigstrut}
\usepackage{multirow}
\usepackage{color}
\usepackage{fancyhdr}
\usepackage{listings} %if lstlistings is used
\usepackage{changepage}
\usepackage{lipsum}
\usepackage{floatrow}
\usepackage{balance}
\usepackage{enumitem}
\usepackage{comment}
\usepackage{breqn}
\usepackage{dblfloatfix}
\usepackage{booktabs} %
\def\BibTeX{{\rm B\kern-.05em{\sc i\kern-.025em b}\kern-.08em
    T\kern-.1667em\lower.7ex\hbox{E}\kern-.125emX}}

\makeatletter
\long\def\@IEEEtitleabstractindextextbox#1{\parbox{0.922\textwidth}{#1}}
\makeatother


%\externaldocument{kiran-raja-ocular-si-revision-1-200521} %[kiran-raja-ocular-si-revision-1-200521.pdf]
%%\externaldocument{kiran-raja-ocular-si-revision-1-200521}

\usepackage{xcite}

\makeatletter
\newcommand*{\addFileDependency}[1]{% argument=file name and extension
	\typeout{(#1)}% latexmk will find this if $recorder=0 (however, in that case, it will ignore #1 if it is a .aux or .pdf file etc and it exists! if it doesn't exist, it will appear in the list of dependents regardless)
	\@addtofilelist{#1}% if you want it to appear in \listfiles, not really necessary and latexmk doesn't use this
	\IfFileExists{#1}{}{\typeout{No file #1.}}% latexmk will find this message if #1 doesn't exist (yet)
}
\makeatother

\newcommand*{\myexternaldocument}[1]{%
	\externaldocument{#1}%
	\addFileDependency{#1.tex}%
	\addFileDependency{#1.aux}%
}
%%% END HELPER CODE

% put all the external documents here!
% put all the external documents here!
\myexternaldocument{MIPGAN-journal-TBIOM-revision-200831}


\newcommand{\Rev}[1]{\textcolor{blue}{#1}}

%% Set counter for review comments 
\newcounter{mycounter} % create a new counter, called 'mycounter'
% default def'n of '\themycounter' is '\arabic{mycounter}'

%% command to increment 'mycounter' by 1 and to display its value:
\newcommand\showReviewcounter{\stepcounter{mycounter}\themycounter}

\usepackage{float}
\makeatletter
\def\bstctlcite{\@ifnextchar[{\@bstctlcite}{\@bstctlcite[@auxout]}}
\def\@bstctlcite[#1]#2{\@bsphack
  \@for\@citeb:=#2\do{%
    \edef\@citeb{\expandafter\@firstofone\@citeb}%
    \if@filesw\immediate\write\csname #1\endcsname{\string\citation{\@citeb}}\fi}%
  \@esphack}
\makeatother 

\title{List of revisions  - "OpenVeinNet: Robust Open-Set Finger Vein Verification with Dynamic Snake Convolution and Graph Learning"}
\begin{document}
\bstctlcite{IEEEexample:BSTcontrol}
%\blackrule[width=\hsize, height=1pt, depth=0.5ex]
\maketitle
\vspace{-2cm}
%\color{blue}
%\color{blue} 
%%%%%%%%%%%% Text goes Below
%\noindent\hrulefill\\
\hrule
\hrulefill
\hrule
\noindent\hrulefill\\
We would like to express our sincere gratitude to the Editor and the reviewers for considering our manuscript for revision. We greatly appreciate the reviewers’ constructive and insightful feedback, which has been invaluable in improving the clarity, technical soundness, and overall quality of the paper. Guided by these comments, we have carefully revised the manuscript and prepared a detailed, point-by-point response addressing each concern. In the following, we summarise the major revisions made to the paper:

\begin{itemize}
    \item Strengthened the motivation and positioning of the proposed OpenVeinNet framework by explicitly linking each methodological component to the characteristics of finger vein recognition: Dynamic Snake Convolution (DSConv) for curvilinear and tubular vein structures, graph-based modelling for long-range vascular connectivity, and the Centroid Angular Hybrid (CAH) loss for compact and angularly separable open-set embeddings. \textbf{(see Main Paper: Sections 1 and 2)}
    
    \item Clarified the novelty of the proposed framework by explicitly stating that DSConv is adapted from prior work, while the contribution of this manuscript lies in its problem-driven integration with graph-based relational modelling and the CAH loss for open-set finger vein verification. \textbf{(see Main Paper: Section 2)}
    
    \item Redesigned and clarified the evaluation protocol to distinguish between cross-dataset full-subject verification and true enrollment-based unknown-rejection evaluation. The revised manuscript now explicitly defines the enrolled and non-enrolled subject groups, the genuine and impostor score generation process, and the threshold-based decision procedure. \textbf{(see Main Paper: Section 3.2)}
    
    \item Added threshold-dependent biometric operating-point metrics, including TAR@FAR=0.1\%, TAR@FAR=1\%, and TAR@FAR=10\%, together with AUC and EER, under both half-subject unknown-rejection and full-subject verification protocols. \textbf{(see Main Paper: Section 3.3, Tables 5 and 6)}
    
    \item Improved reproducibility by correcting the public repository link and providing the implementation, training and evaluation scripts, and model definitions required to reproduce the reported results. The repository README links the \texttt{final\_runs.zip} checkpoint archive, for which users can request the required password and access. \textbf{(see Main Paper: Section 1)}
    
    \item Clarified the experimental repetition strategy by reporting that learned-method results are computed over four predefined identity-disjoint statistical seeds for each leave-one-dataset-out setting and are presented as mean values with 95\% confidence intervals. The intra-database subject-disjoint development/test partition and corresponding image statistics are also reported. \textbf{(see Main Paper: Section 3.2; Supplementary Material: Section 5)}
    
    \item Expanded the baseline training and fairness description by reporting the training configurations of all learned baselines, including optimizer, learning rate, number of epochs, batch size, input resolution, and loss function, under a unified leave-one-dataset-out evaluation protocol. \textbf{(see Supplementary Material: Sections 3 and 4)}
    
    \item Strengthened the ablation study to justify the main design choices of OpenVeinNet, including DSConv kernel sizes, the number of Grapher Blocks, comparison with standard losses, comparison with representative angular-margin and biometric embedding losses, sensitivity analysis of the balancing factor $\beta$, and component-level ablation of DSConv and graph modelling. \textbf{(see Main Paper: Section 4)}
    
    \item Added a dedicated comparison of the proposed CAH loss against representative strong angular-margin and biometric embedding losses, including ArcFace, MagFace, and AdaFace, together with an analysis of the role of the angular balancing factor $\beta$. \textbf{(see Main Paper: Sections 4.4 and 4.5)}
    
    \item Added dataset-quality analysis to explain dataset-specific performance variations, including the strong performance of handcrafted methods on FV-300 and the performance drop on VERA due to low structural coherence, limited samples per identity, and higher acquisition variability. \textbf{(see Main Paper: Section 3.4)}
    
    \item Added computational-cost and deployment analysis, including model size, FLOPs, CPU/GPU latency, and throughput, to clarify the accuracy--efficiency trade-off of the proposed method relative to learned baselines. \textbf{(see Main Paper: Section 5, Table 13)}
    
    \item Added supplementary intra-database open-set verification experiments on FV-300, FV-USM, and MMCBNU to complement the cross-dataset evaluation and provide results under a standard vein-recognition evaluation setting. \textbf{(see Supplementary Material: Section 2)}
    
    \item Added statistical significance testing using paired sign-flip permutation tests with Holm correction to verify that the improvements over learned baselines are statistically significant across matched evaluation runs. \textbf{(see Supplementary Material: Section 1)}
    
    \item Added t-SNE-based embedding visualisations and DET curve analysis in the supplementary material to provide qualitative and operating-characteristic evidence of the learned embedding behaviour across datasets. \textbf{(see Supplementary Material: Sections 6 and 7)}
    
    \item Revised the manuscript language and conclusions to avoid overstatement. The revised manuscript now describes the proposed method as achieving strong overall generalisation and robust open-set verification performance, rather than claiming uniform superiority in every dataset and operating condition. \textbf{(see Main Paper: Sections 3.3, 3.4, and 7)}
\end{itemize}


%------------------------------------------------------------------------------------------------
\color{blue} 
\section*{Co-Editor-in-Chief}
\textit{We have been discussing this paper at the Editorial level. We note that Reviewer 3 noted a lack of clear motivation for the methodological components, suggesting they appear disconnected from the problem. We also note that Reviewer 1 and Reviewer 2 call for substantial revisions. We have gone through the revisions required, but also note the interest expressed by the more supportive reviewers, and even the more negative reviewer's comments on generalisability. As such there is clearly work of interest here, but there appear to be some major improvements to be made consistent with a top journal like TBIOM. These comments suggest that the work should be analysed more carefully in the light of current literature, whilst making several suggestions to improve the paper. Please follow the excellent comments from the reviewers and the AE when preparing the revised version. We shall look closely at the revised version in the light of the comments made.
}
\color{black}
\begin{itemize}
	\item 

We sincerely thank the Co-Editor-in-Chief for the careful editorial assessment and for recognising that the work has potential interest, particularly in relation to generalisation in finger vein verification. We also appreciate the guidance that the revised manuscript should provide stronger motivation, clearer experimental analysis, and better positioning with respect to current literature.

In the revised manuscript, we have substantially reworked the paper to address these concerns. First, we have strengthened the motivation of the proposed methodological components by explicitly linking each part of OpenVeinNet to the characteristics of finger vein images: Dynamic Snake Convolution is motivated by the curvilinear and tubular nature of vein patterns, the graph-based backbone is motivated by the need to model long-range vascular connectivity, and the Centroid Angular Hybrid Loss is motivated by the need for compact and angularly separable embeddings in open-set verification.

Second, we have revised the experimental section to clearly distinguish between cross-dataset generalisation and enrollment-based unknown-rejection evaluation. We now report results under both full-subject verification and stricter open-set unknown-rejection protocols, including AUC, EER, and TAR at fixed FAR operating points. The learned-method results are reported over four predefined identity-disjoint statistical seeds for each leave-one-dataset-out setting, with 95\% confidence intervals to improve statistical reliability.

Third, following the reviewers' suggestions, we have added several new analyses, including stronger angular-margin loss comparisons, beta sensitivity analysis, component-level ablations isolating DSConv and graph modelling, runtime and computational-cost evaluation, dataset-quality analysis, intra-database open-set experiments, statistical significance testing, and t-SNE visualisation of the learned embeddings in the supplementary material.

Finally, we have revised the manuscript language and claims to avoid overstatement. In particular, we now describe the proposed method as achieving strong overall generalisation and robust open-set verification performance, rather than claiming uniform dominance over every method in every dataset condition. We believe these revisions substantially improve the technical clarity, experimental rigour, and suitability of the manuscript for IEEE T-BIOM.

\end{itemize}
%----------------------------------------------------------------------------------------------


%------------------------------------------------------------------------------------------------
\color{blue} 
\section*{Associate Editor}
\textit{ According to the review outcome, the reviewers have strong reservation regarding the manuscript. Particularly, they raised critical concern in the formulation, experimental design and the conclusion drawn. Although the topic appears interesting, it is apparent that more than a major revision is necessary. Based on this observation we regret that the manuscript cannot be accepted. While the outcome may be disappointing to you, we hope you find the reviewers' comments useful for further work in this topic.
}
\color{black}
\begin{itemize}
	\item 

We sincerely thank the Associate Editor for the detailed assessment and for highlighting the key concerns regarding formulation, experimental design, and the conclusions drawn in the original submission. We fully acknowledge that the previous version of the manuscript required substantial improvement in both clarity and technical rigour.

In response, we have thoroughly revised the manuscript to address these concerns. First, the methodological formulation has been significantly strengthened by explicitly motivating each component of OpenVeinNet in relation to the characteristics of finger vein patterns and open-set verification requirements. The roles of Dynamic Snake Convolution, graph-based feature modelling, and the Centroid Angular Hybrid Loss are now clearly justified and cohesively integrated.

Second, the experimental design has been substantially improved. We now provide a well-defined evaluation protocol, including both cross-dataset generalisation and intra-dataset open-set verification settings. The experiments are conducted using identity-disjoint predefined statistical splits, with four learned-method runs per leave-one-dataset-out setting and consistent reporting of AUC, EER, and TAR at fixed FAR operating points. In addition, we have included new analyses such as intra-database experiments, dataset-quality analysis, statistical significance testing, computational cost evaluation, and t-SNE-based embedding visualisation to strengthen the empirical validation.

Third, the conclusions have been revised to better reflect the experimental evidence. We have removed over-generalised claims and now present a more precise interpretation of the results, emphasising strong generalisation and robustness rather than uniform superiority across all scenarios.

We believe that these extensive revisions directly address the concerns raised by the reviewers and significantly improve the clarity, completeness, and scientific contribution of the manuscript.

\end{itemize}
%----------------------------------------------------------------------------------------------
%------------------------------------------------------------------------------------------------
\color{blue} 
%------------------------------------------------------------------------------------------------
\color{blue} 
\section*{Reviewer 1:}
\textbf{\textit{Comment no. \showReviewcounter:}} 
\textit{The first but most important one is the release of the source code and pre-trained model. The link https://github.com/Blazkowiz47/OpenVeinNet does not exist. Hence, the reproducibility is not ensured and the results cannot be verified. This definitely needs to be fixed prior to the publication of the paper.}
\color{black}

\begin{itemize}
	\item \textbf{Authors' Response:} 
	We sincerely thank the reviewer for highlighting the importance of reproducibility and for pointing out the issue with the previously provided repository link. We fully agree that ensuring accessibility of the implementation and trained models is essential for verification and further research. The earlier link was incorrect and resulted from a repository restructuring during development. We apologise for the inconvenience caused.

	\item \textbf{Authors' Action:} 
	We have now made the implementation of the proposed method publicly available at the following repository:
	\url{https://github.com/Blazkowiz47/deep-vein-gcn}. 
	The repository includes the training and evaluation pipeline, model definitions, and scripts required to reproduce the reported results. The repository README links the \texttt{final\_runs.zip} checkpoint archive, and users can request the required password and access for the archive. We have also updated the manuscript to include the correct repository link.
\end{itemize}
\color{blue} 
%%%&&&&&&
%%%&&&&&&
\textbf{\textit{Comment no. \showReviewcounter:}} 
\textit{The authors state that each experiment is repeated with multiple random seeds, but the number of runs is not given.}
\color{black}

\begin{itemize}
	\item \textbf{Authors' Response:} 
	We thank the reviewer for pointing out the need for a clearer specification of the number of experimental runs. We agree that explicitly stating the number of repetitions is important for reproducibility and for correctly interpreting the reported confidence intervals.

	\item \textbf{Authors' Action:} 
	We have clarified this point in the revised manuscript \textbf{(please refer Section 3.2, last paragraph; Page 8)}. The reported learned-method results are obtained using four independent predefined statistical seeds for each leave-one-dataset-out setting. Each seed corresponds to a fixed identity partition stored under the corresponding \texttt{data/<dataset>/<stat\_seed>} directory, resulting in four trained models and four corresponding evaluations for each held-out test dataset.

	The final performance is reported as the {mean and 95\% confidence interval over these four runs} for each leave-one-dataset-out protocol. For the paired significance analysis, the matched observations are pooled across the five held-out datasets and four statistical seeds, giving \(n=20\) matched runs. This clarification has been explicitly added to the experimental protocol section to improve transparency and reproducibility.
\end{itemize}
\color{blue} 
%%%&&&&&&
%%%&&&&&&
%%%&&&&&&
\textbf{\textit{Comment no. \showReviewcounter:}} 
\textit{The text for Table 3 states that some “bad quality” samples were removed from the FV-300 for the experiments. There is no information on which or how many samples were removed.}
\color{black}

\begin{itemize}
	\item \textbf{Authors' Response:} 
	We thank the reviewer for highlighting the need for clearer reporting of the data preprocessing step. We agree that explicitly specifying both the number of removed samples and the criteria used for removal is important for transparency and reproducibility.

	\item \textbf{Authors' Action:} 
	We have clarified this point in the revised manuscript \textbf{(please refer Section 3.1, first paragraph, last sentence and Table 3 Caption in  Page 7)}. Specifically, we now explicitly state that {124 images were removed from the FV-300 dataset}. These samples were excluded due to {severe blur, poor illumination, or failure in vein visibility}, which make them unsuitable for reliable feature extraction and verification.

	The corresponding clarification has been added to Table~3 and its caption, as well as in the dataset description, to ensure that both the number of removed samples and the filtering criteria are clearly documented.
\end{itemize}

\color{blue} 
%%%&&&&&&
\textbf{\textit{Comment no. \showReviewcounter:}} 
\textit{The authors state that the samples per identity are deliberately varied between training and validation sets, but the training and validation split is not given.}
\color{black}

\begin{itemize}
	\item \textbf{Authors' Response:} 
	We thank the reviewer for pointing out this missing protocol detail. We agree that the partitioning strategy should be explicitly reported, as it affects reproducibility and the interpretation of the training protocol.

	In the revised manuscript and supplementary material, we clarify that the intra-database train--test split is performed at the \textbf{subject/identity level}, not by randomly splitting images within the same identity. For each statistical split, identities are partitioned into disjoint development and held-out test subsets. Images belonging to development identities are used only for model development, while images belonging to held-out test identities are used only for final testing. Thus, no subject or finger identity used for final open-set testing appears during model development.

	The image counts reported in the supplementary material are obtained after this subject/identity-disjoint partitioning by counting the images associated with the identities assigned to each subset. These counts are included to clarify the amount of data used for development and held-out testing, as requested by the reviewer, but they do not imply image-level splitting within the same identity.

	\item \textbf{Authors' Action:} 
	We have clarified the train--test partitioning strategy in the revised manuscript \textbf{(please refer Section 3.2, last paragraph; Page 8)} and provide the corresponding per-dataset train--test statistics in the supplementary material \textbf{(please refer Section 5 of the Supplementary Material; Page 5)}. The revised description explicitly states that:
	\begin{itemize}
		\item the split is performed at the subject/identity level;
		\item no subject or finger identity is shared between model development and held-out testing;
		\item all images belonging to a development identity are used only for model development;
		\item all images belonging to a held-out test identity are used only for final testing; and
		\item the reported image counts are obtained after subject/identity-disjoint partitioning.
	\end{itemize}

	This clarification ensures that the training and testing protocol is transparent, reproducible, and free from identity leakage.
\end{itemize}
%%%&&&&&&
\color{blue} 
\textbf{\textit{Comment no. \showReviewcounter:}} 
\textit{The authors state that their approach outperforms all the tested approaches from the literature. However, there are some approaches specifically tailored to open-set vein recognition, e.g. [1], which have not been tested by the authors. Furthermore, no details about the training and parameter optimisation of the tested approaches are provided in the paper. }
\color{black}

\begin{itemize}
	\item \textbf{Authors' Response:} 
	We sincerely thank the reviewer for highlighting this important point regarding the inclusion of open-set-specific methods and the clarity of experimental settings. We fully agree that comparison with methods explicitly designed for open-set vein recognition is highly relevant.

	In the revised manuscript, we have carefully incorporated the work of Chen \emph{et al.} (2021) [18] as a representative open-set finger vein recognition approach. In addition, we have revisited our experimental protocols to ensure closer alignment with the evaluation principles adopted in that work, particularly with respect to open-set verification and cross-dataset generalisation.

	However, we regret that a direct quantitative comparison with Chen \emph{et al.} [18] could not be reliably established. Specifically, (i) the official implementation and trained models are not publicly available, and (ii) our independent re-implementation, based on the methodological description provided in the paper, did not reproduce results consistent with those reported. We believe this discrepancy arises from missing details related to training strategy, data preprocessing, and hyperparameter configuration, which are not fully specified in the original work.

	We emphasise that including results from a non-validated re-implementation could lead to unfair or misleading comparisons. Therefore, in the interest of maintaining scientific rigour and reproducibility, we have opted not to report these results quantitatively. Instead, we explicitly acknowledge this limitation and position our evaluation against a broad set of both handcrafted and deep learning-based baselines under a unified and reproducible experimental framework.

	Regarding the reviewer’s second point, we agree that the description of training and parameter optimisation should be clearer. We have now expanded the manuscript to provide additional details on the training setup, including optimisation strategy, hyperparameter choices, and evaluation protocol, ensuring that all compared methods are implemented under consistent and fair conditions.

	\item \textbf{Authors' Action:} 
	\begin{itemize}
		\item Included Chen et al. (2021) as an open-set reference (Ref. [18]) in the related work discussion and literature summary table (please refer Table 1, Page 2 in main paper).
		\item Aligned the experimental protocols  with open-set evaluation practices described in prior work. \textbf{(please refer Section 3.2, page 7, in main paper})
		\item Expanded the results section to reflect the new results based on open-set evaluation practices (\textbf{please refer Table 5 and 6, Page 9 and 10, in main paper}).
        \item Included a new section (\textbf{please refer Supplementary Material Section 4}) to provide clearer descriptions of training procedures, parameter settings, and evaluation protocols for all compared methods.
	\end{itemize}
\end{itemize}
%%%&&&&&&
%%%&&&&&& 
\color{blue} 
\textbf{\textit{Comment no. \showReviewcounter:}} \textit{For the ablation study it is not clear which datasets or which open-set scenario has been used. Only for the loss function ablation study the authors state that it was ABDE -> C, but for the number of Grapher blocks and kernel sizes, this information is missing. }
\color{black}

\begin{itemize}
\item \textbf{Authors' Response:} 
We thank the reviewer for pointing out this lack of clarity. We agree that the description of the experimental protocol used in the ablation study was not consistently stated across all subsections. While the loss ablation explicitly mentioned the evaluation setting, the same protocol was implicitly used for the other ablation experiments (kernel-size configuration and Grapher Block analysis), which may have caused confusion.

To clarify, all ablation experiments were conducted under the same open-set protocol, namely $ABDE \rightarrow C$, where FV-USM is used as the held-out evaluation dataset. This ensures that the comparisons across architectural and loss components are consistent and directly comparable.

\item \textbf{Authors' Action:} 
We have revised Section~4 to explicitly state that all ablation experiments are performed under the $ABDE \rightarrow C$ protocol. In addition, we added a unifying statement at the beginning of the ablation section to clearly indicate the evaluation setting used throughout the analysis. (\textbf{please refer Section 4, first paragraph, in main paper})

\end{itemize}
\color{blue} 
%%%&&&&&&
%%%&&&&&&
\textbf{\textit{Comment no. \showReviewcounter:}} 
\textit{The text for Fig. 6 refers to Block 1 and Pooling, but in the figure there is neither a Block 1 nor a pooling layer.}
\color{black}

\begin{itemize}
	\item \textbf{Authors' Response:} 
	We thank the reviewer for carefully pointing out this inconsistency between the figure and the accompanying text. We agree that the previous version of Fig.~6 did not clearly show all the components referred to in the manuscript, which could make the backbone description difficult to follow.

	\item \textbf{Authors' Action:} 
	We have revised Fig.~6 to explicitly include the backbone components described in the text, including \textbf{Block 0}, \textbf{Block 1}, and the \textbf{adaptive average pooling} layer. We have also revised the corresponding text and figure caption to ensure consistency between the visual architecture and the written description.
\end{itemize}
%%%&&&&&&

%%%&&&&&&
\color{blue} 
%%%&&&&&&
\textbf{\textit{Comment no. \showReviewcounter:}} 
\textit{The Fig. 7 plots are too small and hard to read.}
\color{black}

\begin{itemize}
	\item \textbf{Authors' Response:} 
	We thank the reviewer for highlighting the readability issue in Fig.~7. We agree that the previous version of the plots was difficult to interpret due to their small size and dense visual content.

	\item \textbf{Authors' Action:} 
	In the revised manuscript, we have removed the DET curve plots from the main paper to improve clarity and maintain focus on the quantitative evaluation results presented in Tables~5 and~6. These tables already report the key operating-point metrics (AUC, EER, and TAR@FAR), which provide a clearer and more compact comparison across methods.

	To ensure completeness and reproducibility, we have included the DET curve visualisations in the supplementary material, where they are presented at higher resolution and analysed without space constraints.
\end{itemize}
%%%&&&&&&



%------------------------------------------------------------------------------------------------
\setcounter{mycounter}{0}
\color{blue} 
\section*{Reviewer  2:}
%%%&&&&&&
\textbf{\textit{Comment no. \showReviewcounter:}} 
\textit{The paper asserts to address ``open-set verification,'' yet the described evaluation methodology seems to be cross-dataset evaluation rather than true open-set verification. In true open-set verification, the system must not only handle unseen subjects but also correctly reject impostors who do not belong to any enrolled identity. The current evaluation protocol (training on four datasets and testing on the fifth) assesses cross-domain generalization but does not explicitly test the system's ability to reject unknown identities. It is necessary to clarify whether the evaluation truly constitutes open-set verification or whether it would be more precisely described as cross-dataset closed-set verification. If it is the former, additional experiments demonstrating false acceptance rates at specific false rejection rates would strengthen this claim.}
\color{black}

\begin{itemize}
	\item \textbf{Authors' Response:} 
	We sincerely thank the reviewer for this important clarification. We agree that the original manuscript did not sufficiently distinguish between cross-dataset subject-disjoint verification and true enrollment-based open-set verification. The reviewer is correct that training on four datasets and testing on a held-out dataset alone mainly evaluates cross-dataset generalisation, unless the evaluation explicitly includes non-enrolled identities that must be rejected.

	To address this concern, we have redesigned and clarified the evaluation protocol in the revised manuscript. We now report two complementary protocols. The first is an \textbf{enrollment-based unknown-rejection protocol}, where only 50\% of the subjects from the held-out test dataset are enrolled, while the remaining 50\% are treated as unknown identities. Genuine scores are computed from all within-identity comparisons among enrolled subjects, and impostor scores are computed exhaustively from cross-identity comparisons between enrolled and non-enrolled subjects. FAR, FRR, EER, AUC, and TAR@FAR are then derived from these pairwise score distributions by sweeping the decision threshold. This explicitly evaluates whether enrolled-subject scores can be separated from non-enrolled unknown-subject scores. The second protocol is a \textbf{full-subject verification protocol}, where all subjects from the held-out test dataset are included in the pairwise verification evaluation, corresponding to subject-disjoint cross-dataset verification.

	In addition, following the reviewer’s suggestion, we now report threshold-dependent operating points, including TAR@FAR$=0.1\%$, TAR@FAR$=1\%$, and TAR@FAR$=10\%$, together with AUC and EER. These operating points directly quantify the verification behaviour at fixed false-acceptance regimes and therefore provide a more deployment-relevant assessment of open-set performance.

	\item \textbf{Authors' Action:} 
	We have revised \textbf{Section~3.2, in main paper} to clearly define the two evaluation protocols and to distinguish enrollment-based unknown-rejection from full-subject verification. We have also added new result tables in Section~3.3 reporting AUC, EER, and TAR@FAR operating points for both protocols. The revised results demonstrate that the proposed method performs strongly not only under cross-dataset verification, but also under the stricter open-set unknown-rejection setting requested by the reviewer.
\end{itemize}
%%%&&&&&&
%%------------------
\color{blue} 

%%%&&&&&&
\textbf{\textit{Comment no. \showReviewcounter:}} 
\textit{The paper compares with seven state-of-the-art methods; however, several concerns regarding the fairness of these comparisons emerge:
MCP, RLT, and WLD are traditional hand-crafted methods, and it is surprising that they outperform many deep learning methods (e.g., ArcVein, LGFIN, FV-ViT) on certain datasets. This raises questions about whether the deep learning baselines were properly implemented and optimized.
Table 5 shows that MCP achieves 99.80\% AUC on BCDE$\rightarrow$A, while deep learning methods like ArcVein achieve only 89.77\%. This significant difference is unusual and requires an explanation—are the deep learning baselines using the same evaluation protocol?
The paper should clarify whether the baseline results were reproduced using the original authors' implementations or taken from published literature and ensure that all methods use identical training/testing splits.}
\color{black}

\begin{itemize}
	\item \textbf{Authors' Response:} 
	We thank the reviewer for this important observation regarding the fairness of the comparisons. We agree that the strong performance of handcrafted methods (MCP, RLT, WLD) in certain settings, particularly the BCDE$\rightarrow$A protocol, requires careful explanation, and that the training and evaluation setup of the deep baselines must be clearly stated.

	First, we clarify that {all learned baselines reported in the revised manuscript (ArcVein, LGFIN, FV-ViT, and VeinAttNet) were re-trained by us} under a unified experimental framework, rather than directly taken from the literature. Specifically, all learned methods were trained using the same leave-one-dataset-out protocol, the same subject-disjoint splits, matched statistical seeds, and a common evaluation pipeline for AUC, EER, and TAR computation. This ensures that the reported results are directly comparable across methods.

	Second, the observation that handcrafted methods outperform several deep models in the BCDE$\rightarrow$A setting (i.e., FV-300 as the target dataset) is consistent with the properties of the dataset. As shown in the newly added \emph{Dataset Quality Analysis} (Section~3.4), FV-300 is a relatively clean and well-controlled dataset with high image quality and a large number of samples per identity. In such conditions, handcrafted vessel-enhancement and template-based matching methods remain highly competitive, since the vein patterns are already sharp and well-aligned. In contrast, deep learning models trained under cross-dataset conditions may experience a domain shift that slightly degrades performance when transferred to this specific dataset.

	Therefore, this behaviour does not indicate an unfair comparison or improper optimisation of the deep baselines, but rather reflects a known characteristic of finger-vein recognition: handcrafted methods can perform strongly on clean, high-quality datasets, while deep models show their advantage primarily in cross-dataset and low-quality scenarios. This is further supported by the results on more challenging datasets (e.g., FV-USM, MMCBNU, and VERA), where the proposed method consistently outperforms both handcrafted and deep baselines.

	\item \textbf{Authors' Action:} 
	To address this concern, we have made the following revisions:
	\begin{itemize}
	    \item Added a detailed description of the {baseline training configurations and fairness protocol} in the supplementary material, including optimizers, learning rates, batch sizes, and training schedules for all learned methods. \textbf{(please refer Section 4 in supplementary material)}
	    \item Explicitly clarified in the manuscript that {all deep baselines were re-trained under the same experimental protocol}, rather than using numbers from the literature. \textbf{(please refer Section 4 in supplementary material)}
	    \item Introduced a new subsection on {Dataset Quality Analysis and Performance Interpretation} to explain the strong performance of handcrafted methods on FV-300 and to provide a principled interpretation of the observed results. \textbf{(please refer Section 3.4 in main paper.)}
	\end{itemize}
	These additions ensure that the comparisons are transparent, fair, and technically justified.
\end{itemize}
%%%&&&&&&

%%------------------
\color{blue} 
%%%&&&&&&
\textbf{\textit{Comment no. \showReviewcounter:}} 
\textit{Although the paper proposes several novel elements, the claims of contribution need stronger justification:
Dynamic Snake Convolution: The paper states that DSConv "adjusts its sampling locations in a flexible manner that follows the natural contours of the vein patterns." However, DSConv appears to be adapted from existing work (possibly from medical imaging applications). It is necessary to clarify the novelty of DSConv in the context of finger vein recognition and cite the original DSConv work.
Centroid-Angular Hybrid Loss: This loss function combines softmax with an angular component. The mathematical formulation in Equations 2--3 is intricate and would benefit from: (a) a clearer explanation of the intuition behind the angular term, (b) ablation studies demonstrating the effect of the $\beta$ parameter, and (c) comparison with existing margin-based losses (ArcFace, CosFace) used in face recognition. The paper mentions releasing source code, which is beneficial for reproducibility. Please ensure that the code is properly documented and includes training configurations for all experiments.}
\color{black}

\begin{itemize}
	\item \textbf{Authors' Response:} 
	We sincerely thank the reviewer for this constructive comment. We agree that the original manuscript did not sufficiently separate the novelty of adopting DSConv for finger vein verification from the original DSConv formulation, nor did it provide enough intuition and empirical support for the proposed Centroid-Angular Hybrid (CAH) loss.

	Regarding DSConv, we have revised the manuscript to clearly state that DSConv itself is adapted from prior work and have cited the original DSConv paper. Our contribution is not the invention of DSConv, but its integration into an open-set finger vein verification framework, where it is used as a task-specific stem to capture tubular and curvilinear vascular patterns before graph-based relational modelling. This clarification avoids overstating the novelty and more accurately positions the contribution in the context of finger vein recognition.

	Regarding the CAH loss, we have expanded the explanation of the angular component. The revised text now explains that the angular term is intended to encourage identity embeddings to align with class-specific direction vectors while maintaining compact intra-class structure and improved inter-class separation. This is particularly important for open-set verification, where unknown identities must be separated from enrolled identities in the embedding space.

	We have also added new ablation studies to directly address the reviewer’s requests. First, we report a sensitivity analysis of the balancing factor $\beta$, showing that $\beta=0.3$ provides the best trade-off between classification supervision and angular regularisation. Second, we compare the proposed CAH loss against representative strong angular-margin and biometric embedding losses, including ArcFace, MagFace, and AdaFace, under the same ABDE $\rightarrow$ C protocol. These losses were selected because ArcFace is a widely used angular-margin baseline, while MagFace and AdaFace are more recent strong biometric representation losses that account for feature quality and adaptive margins.

	Finally, we have improved the reproducibility material by providing the public code repository and adding detailed training configurations for all experiments and baselines in the supplementary material.

	\item \textbf{Authors' Action:} 
	We have made the following revisions:
	\begin{itemize}
	    \item Added the original DSConv citation and clarified that DSConv is adapted from prior work, while our novelty lies in its use within an open-set finger vein verification framework together with graph-based modelling and CAH loss. \textbf{(please refer Section 2, para 1, citation number 43 in the main paper.)}
	    \item Revised the proposed method section to better explain the intuition behind the angular term in the CAH loss. \textbf{(please refer Section 2.3 in the main paper.)}
	    \item Added a $\beta$ sensitivity study in the ablation section. \textbf{(please refer Section 4.5 in the main paper.)}
	    \item Added a comparison against representative strong angular-margin and biometric embedding losses, including ArcFace, MagFace, and AdaFace. (please refer Section 4.4 in the main paper.)
	    \item Updated the public repository link and added detailed baseline training configurations in the supplementary material to support reproducibility. \textbf{(please refer Section 4 in the Supplementary Material.)}
	\end{itemize}
\end{itemize}
%%%&&&&&&
%%------------------
\color{blue} 
%%%&&&&&&
\textbf{\textit{Comment no. \showReviewcounter:}} 
\textit{Architectural Complexity and Computational Cost
The proposed architecture (approximately 4.7 million parameters) is relatively complex, consisting of multiple components (DSConv stem, Grapher blocks, hybrid loss). For practical deployment, especially in resource-constrained environments, the following should be provided: Inference time (FPS) on GPU and CPU, comparison of model size with baseline methods, and a discussion of the computational trade-offs versus accuracy gains.}
\color{black}

\begin{itemize}
	\item \textbf{Authors' Response:} 
	We thank the reviewer for this important suggestion regarding deployment considerations. We agree that, in addition to accuracy, computational efficiency is a key factor for practical biometric systems.

	In the revised manuscript, we have included a dedicated analysis of computational cost and inference latency. Specifically, we now report the number of parameters, FLOPs, inference latency (ms), and throughput (images per second) on both CPU and GPU for all learned methods under a unified experimental setting. This allows a direct comparison between the proposed method and existing baselines.

	The results show that while lightweight models such as FV-ViT and ArcVein achieve lower latency and higher throughput, they exhibit reduced verification performance, particularly under challenging cross-dataset and open-set conditions. In contrast, the proposed OpenVeinNet has a moderate model size (4.76M parameters) and computational cost (5.07 GFLOPs), with higher latency (55.82 ms CPU / 13.00 ms GPU), but consistently delivers superior accuracy across all evaluation protocols.

	From a design perspective, the increased computational cost is primarily due to the DSConv layers and Grapher blocks. DSConv introduces adaptive sampling aligned with curvilinear vein structures, while the Grapher blocks enable dynamic graph construction and long-range relational modelling. These components are inherently more computationally intensive than standard convolutions but directly contribute to improved discriminative capability and generalisation.

	Overall, the results highlight a clear accuracy--efficiency trade-off. The proposed method prioritises robustness and verification performance, while maintaining a computational footprint that remains practical for deployment in GPU-enabled or server-side environments.

\item \textbf{Authors' Action:} 
We have added a new section including:
\begin{itemize}
    \item A detailed table reporting model size, FLOPs, CPU/GPU latency, and throughput for all learned methods. \textbf{(please refer Section 5, Table 13 in the main paper.)}
    \item A discussion of the accuracy--efficiency trade-off and the computational cost introduced by DSConv and Grapher Blocks. \textbf{(please refer Section 5 in the main paper.)}
\end{itemize}
This addition provides a comprehensive and transparent evaluation of the proposed method from a deployment perspective.
\end{itemize}
%%%&&&&&&


\color{blue} 
%%------------------
%%%&&&&&&
\textbf{\textit{Comment no. \showReviewcounter:}} 
\textit{Generalization Analysis
The cross-dataset evaluation is rigorous and praiseworthy. However, the paper could strengthen its analysis by: providing statistical significance tests to determine whether performance differences between OpenVeinNet and the second-best method are statistically significant; discussing the reasons for the better generalization of the method, perhaps through additional analysis of the learned features or visualization of the feature space; and addressing the performance drop on the VERA dataset (89.98\% AUC versus >97\% on other datasets), which suggests vulnerability to certain types of noise or acquisition conditions.}
\color{black}

\begin{itemize}
	\item \textbf{Authors' Response:} 
	We sincerely thank the reviewer for this insightful comment. We agree that a deeper analysis of generalisation behaviour is important to support the empirical results.

	First, regarding statistical significance, we have added formal pairwise significance testing between the proposed method and the competing learned baselines. Specifically, we performed two-sided exact sign-flip permutation tests on the full-subject EER results across matched leave-one-dataset-out runs, with Holm correction applied for multiple comparisons. The results show that the improvements of OpenVeinNet over the compared learned baselines are statistically significant, confirming that the observed performance gains are consistent rather than due to random variation.

	Second, to better understand the generalisation capability of the proposed method, we have added feature-space analysis in the supplementary material using t-SNE visualisations of the OpenVeinNet embeddings. These visualisations show that the proposed model produces meaningful intra-class grouping and dataset-dependent separability patterns. In particular, datasets with higher image quality and more samples per identity show more compact clusters, whereas low-sample and high-variability datasets such as VERA show more dispersed embeddings. These observations support the quantitative trends reported in the main paper.

	Third, regarding the performance drop on the VERA dataset, we have added a dedicated analysis in the revised manuscript. The results show that VERA is the most challenging dataset due to two key factors: (i) it has the lowest gradient-coherence score among all datasets, indicating weaker structural continuity of vein patterns, and (ii) it contains only two images per finger identity/class in the evaluation setting, resulting in very limited intra-class variation during evaluation. This combination leads to a low-sample, high-variability test condition, which naturally degrades performance for all methods. Therefore, the observed drop is not specific to the proposed method, but reflects the inherent difficulty of the dataset.

	\item \textbf{Authors' Action:} 
	We have made the following additions to the manuscript and supplementary material:
	\begin{itemize}
	    \item Included statistical significance testing results, with permutation tests and Holm correction, in the supplementary material. \textbf{(please refer Section 1, Table 1 in the Supplementary Material.)}
	    \item Added t-SNE-based feature-space visualisations and corresponding discussion to analyse embedding separability and generalisation behaviour. \textbf{(please refer Section 6 in the Supplementary Material.)}
	    \item Introduced a new subsection on \textit{Dataset Quality Analysis and Performance Interpretation} to explain the VERA performance drop and the behaviour of different methods across datasets. \textbf{(please refer Section 3.4 in the main paper.)}
	\end{itemize}
\end{itemize}
%%%&&&&&&
%------------------------------------------------------------------------------------------------
\setcounter{mycounter}{0}
\color{blue} 
\section*{Reviewer  3:}

%%------------------
\color{blue}
\textbf{\textit{Comment no. \showReviewcounter:}} 
\textit{The proposed framework essentially combines two existing components: snake convolution for vascular feature extraction and GNN-based relational modeling. The integration appears straightforward, with very limited insights (analytically or experimentally) to the field of finger vein recognition.}
\color{black}

\begin{itemize}
\item \textbf{Authors' Response:} 
We sincerely thank the reviewer for this thoughtful and important observation. We agree that the individual components used in our framework, dynamic snake convolution and graph-based modelling, are not entirely novel in isolation. However, our contribution lies in the \emph{problem-driven integration} of these components within an open-set finger vein verification setting, along with the design of a dedicated embedding objective.

Specifically, the motivation of the proposed architecture is rooted in the structural properties of finger vein patterns. Finger veins exhibit elongated, curvilinear, and spatially connected structures. Standard convolutional layers are limited in capturing such geometry, while graph-based methods alone lack the ability to extract robust local tubular features. In this work, DSConv is used as a \emph{task-specific stem} to adaptively align convolutional sampling with vein structures, while the Grapher backbone explicitly models long-range topological relationships between these structures. This combination is therefore not a straightforward stacking of modules, but a \emph{structurally complementary design} tailored to the finger vein domain.

In addition, the proposed Centroid-Angular Hybrid (CAH) loss is introduced to address open-set verification requirements by enforcing both intra-class compactness and angular separation between identities. This is particularly important in cross-dataset and unknown-rejection scenarios, where the embedding space must generalise beyond seen identities.

Importantly, we have strengthened the manuscript with \emph{extensive experimental evidence} to support these design choices. The ablation study shows that:
(i) DSConv alone significantly improves performance over a standard convolutional stem,
(ii) graph-based modelling independently contributes to performance gains, and
(iii) the combination of both yields the best results, demonstrating a clear synergistic effect rather than a trivial combination.
Furthermore, the loss ablation shows that CAH outperforms both standard and strong angular-margin losses.

\item \textbf{Authors' Action:} 
To better reflect the contribution and address the reviewer’s concern, we have:
\begin{itemize}
    \item Revised the introduction and method sections to clearly articulate the {problem-driven motivation} behind combining DSConv and graph-based modelling for finger vein structures.  \textbf{(please refer Section 1 and 2  in the main paper.)}
    \item Explicitly clarified that DSConv is adapted from prior work, while the novelty lies in its {integration with graph modelling and CAH loss} for open-set verification. \textbf{(please refer Section  2  in the main paper.)}
    \item Strengthened the {ablation study} to highlight the individual and combined contributions of DSConv, Grapher blocks, and the proposed loss. \textbf{(please refer Section 4  in the main paper.)}
    \item Added additional analysis, including feature-space visualisation in the supplementary material and dataset-quality discussion in the main paper, to provide deeper insights into the model behaviour. (please refer Section 6 in the Supplementary Material and Section 3.4 in the main paper.)
\end{itemize}

These revisions aim to better communicate the technical insight and experimental justification underlying the proposed framework.
\end{itemize}
%%%&&&&&&
%%------------------
\color{blue} 
%%%&&&&&&
\textbf{\textit{Comment no. \showReviewcounter:}} 
\textit{The proposed loss function is presented as a reformulation of the softmax loss. However, similar angular-margin-based losses (e.g., ArcFace, CosFace, SphereFace) have been extensively studied and widely adopted in biometric recognition tasks. The manuscript does not clearly explain the motivation behind the proposed loss formulation, nor the theoretical or practical advantages over existing angular-margin losses. Without clear justification or comparison with state-of-the-art losses, the novelty and significance of the proposed loss remain unclear.}
\color{black}

\begin{itemize}
	\item \textbf{Authors' Response:} 
	We sincerely thank the reviewer for this important and insightful comment. We agree that the original manuscript did not sufficiently position the proposed Centroid-Angular Hybrid (CAH) loss with respect to existing angular-margin losses such as ArcFace, CosFace, and SphereFace.

	In the revised manuscript, we clarify that the proposed loss is not intended as a direct replacement of margin-based losses, but rather as a \emph{hybrid objective} tailored for open-set finger vein verification. While existing angular-margin losses typically enforce angular or cosine margins between classes in a classification setting, they are primarily designed for closed-set recognition scenarios. In contrast, the CAH loss combines a centroid/prototype-based formulation with angular regularisation, which encourages both compact intra-class clustering and stable directional structure in the embedding space. This is particularly important in open-set and cross-dataset conditions, where the model must generalise to unseen identities and maintain separability under domain shift.

	Intuitively, the centroid/prototype component stabilises the representation by encouraging samples to align with identity-specific class directions, while the angular term enforces directional consistency in the embedding space. This combination is intended to improve robustness under domain shift and limited-sample regimes, which are common in finger vein datasets.

	Furthermore, we have added a direct comparison with representative strong angular-margin and biometric embedding losses, including ArcFace, MagFace, and AdaFace, under the same ABDE $\rightarrow$ C protocol. These losses were selected because ArcFace is a widely used angular-margin baseline, while MagFace and AdaFace are more recent strong biometric representation losses that further account for feature quality and adaptive margins. The proposed CAH loss achieves the lowest EER among these compared objectives, supporting its effectiveness for the proposed open-set finger vein verification setting. We also added a sensitivity analysis of the balancing factor $\beta$, showing how the angular component contributes to performance when properly weighted.

	\item \textbf{Authors' Action:} 
	To address this concern, we have made the following revisions:
	\begin{itemize}
	    \item Expanded the description of the proposed loss to clearly explain the \emph{motivation and intuition} behind the centroid/prototype and angular components. \textbf{(please refer Section 2.3 in the main paper.)}
	    \item Added a comparison against representative strong angular-margin and biometric embedding losses, including ArcFace, MagFace, and AdaFace. \textbf{(please refer Section 4.4 in the main paper.)}
	    \item Included a \emph{sensitivity analysis of the $\beta$ parameter} to demonstrate the role of the angular term. \textbf{(please refer Section 4.5 in the main paper.)}
	    \item Revised the text to explicitly position the CAH loss as a \emph{task-specific hybrid formulation} for open-set verification, rather than a generic replacement of existing margin-based losses. \textbf{(please refer Section 2.3 in the main paper.)}
	\end{itemize}

	These additions clarify both the design rationale and the practical advantages of the proposed loss in the context of finger vein verification.
\end{itemize}
%%%&&&&&&

%%------------------
\color{blue} 
%%%&&&&&&
\textbf{\textit{Comment no. \showReviewcounter:}} 
\textit{The manuscript repeatedly claims that the proposed method outperforms both handcrafted and deep learning approaches across five datasets. However, the results reported in Table 5 contradict this statement. Specifically, the proposed method performs significantly worse than traditional handcrafted methods such as MCP, RLT, and WLD on certain datasets. This abnormal observation is not discussed or analyzed in the manuscript. Such inconsistencies weaken the credibility of the conclusions.}
\color{black}

\begin{itemize}
	\item \textbf{Authors' Response:} 
	We sincerely thank the reviewer for pointing out this important inconsistency. We agree that the original wording in the manuscript was overly strong and did not accurately reflect the behaviour observed in Table~5. In particular, the performance of handcrafted methods on certain datasets (e.g., FV-300 in the $BCDE \rightarrow A$ setting) was not sufficiently discussed, which could indeed affect the credibility of the conclusions.

	In the revised manuscript, we have carefully analysed this observation and clarified the claims. The strong performance of handcrafted methods such as MCP, RLT, and WLD in specific settings is primarily due to the characteristics of the dataset. As shown in the newly added \emph{Dataset Quality Analysis}, FV-300 is a high-quality and well-controlled dataset with a large number of samples per identity. In such conditions, handcrafted vessel-enhancement and template-based methods remain highly effective, as the vein structures are already clear and consistent, requiring less complex feature learning.

	In contrast, the proposed method is designed to handle more challenging scenarios, including cross-dataset generalisation, low-quality images, and limited samples per identity. This is reflected in the results on more difficult datasets (e.g., FV-USM, MMCBNU, and VERA), where the proposed method achieves strong overall performance and shows clear gains over most handcrafted and deep learning baselines, especially in terms of EER and threshold-based verification behaviour.


	\item \textbf{Authors' Action:} 
	To address this concern, we have made the following revisions:
	\begin{itemize}
	    \item Revised the claims throughout the manuscript to state {strong overall generalisation performance}, rather than implying dominance on every dataset.
	    \item Added a new subsection on {Dataset Quality Analysis and Performance Interpretation} to explicitly explain the strong performance of handcrafted methods on FV-300. \textbf{(please refer Section 3.4 in the main paper.)}
	    \item Updated the discussion around Table~5 (from main paper) to provide a balanced and technically grounded interpretation of the results.
	\end{itemize}

	These changes ensure that the conclusions are accurate, transparent, and aligned with the experimental evidence.
\end{itemize}
%%%&&&&&&
%%------------------

%%%&&&&&&
\color{blue} 
\textbf{\textit{Comment no. \showReviewcounter:}} 
\textit{In Table 4 (Section 3.2), the protocol for constructing genuine and impostor pairs is not clearly described, e.g., how the genuine and impostor pairs are generated, whether all possible pairs are used or sampling is performed? The lack of protocol clarity makes it difficult to assess the validity and reproducibility of the reported results.}
\color{black}

\begin{itemize}
	\item \textbf{Authors' Response:} 
	We thank the reviewer for highlighting this important point regarding the clarity of the evaluation protocol. We agree that the original manuscript did not sufficiently describe how genuine and impostor pairs are constructed, which is essential for assessing the validity and reproducibility of the reported results.

	In the revised manuscript, we have explicitly clarified the pair-generation protocol used for both evaluation settings. For each evaluation split, \textbf{genuine pairs} are generated exhaustively by forming all within-identity image comparisons, excluding self-comparisons. Therefore, every valid pair of images belonging to the same finger identity contributes to the genuine score distribution.

	\textbf{Impostor pairs} are also generated exhaustively, without random sampling. Only cross-identity comparisons are used. The exact definition of impostor pairs depends on the evaluation protocol. In the \textbf{half-subject enrollment-based unknown-rejection protocol}, only 50\% of the subjects in the held-out test dataset are enrolled, while the remaining 50\% are treated as non-enrolled unknown identities. Genuine scores are computed from all within-identity comparisons among enrolled subjects. Impostor scores are computed from all cross-identity comparisons between samples from enrolled subjects and samples from non-enrolled subjects. This directly evaluates whether the system can reject unknown identities that are not present in the enrolled gallery.

	In the \textbf{full-subject verification protocol}, all subjects from the held-out test dataset are included in the evaluation. Genuine scores are computed from all within-identity image comparisons, excluding self-comparisons, and impostor scores are computed from all cross-identity comparisons between different identities in the held-out dataset. Thus, the reported FAR, FRR, EER, AUC, and TAR@FAR values are derived from exhaustive pairwise score distributions rather than from randomly sampled pairs.

	All methods are evaluated using the same protocol and the same genuine/impostor score-generation strategy under each setting. No method-specific sampling, filtering, or pair selection is applied. This ensures that the reported comparisons are fair, transparent, and reproducible.

	\item \textbf{Authors' Action:} 
	We have revised \textbf{Section~3.2 in the main paper} to include a detailed description of the genuine and impostor pair construction process, including:
	\begin{itemize}
	    \item Explicit definition of genuine and impostor pairs.
	    \item Clarification that genuine pairs are generated exhaustively within each identity, excluding self-comparisons.
	    \item Clarification that impostor pairs are also generated exhaustively, without random sampling.
	    \item Explanation of the difference between impostor-pair construction in the half-subject unknown-rejection protocol and the full-subject verification protocol.
	    \item Statement that all methods are evaluated using the same pair-generation strategy for fairness.
	\end{itemize}

	This addition makes the evaluation protocol transparent and fully reproducible.
\end{itemize}
%%%&&&&&&------

\color{blue} 
%%%&&&&&&
\textbf{\textit{Comment no. \showReviewcounter:}} 
\textit{In Table 5 (Section 3.3), the authors report cross-dataset evaluation results, which is indeed a valuable indicator of generalization ability. However, the reported performance under this setting is actually very poor, with EER values ranging from 3\% to 18\% by the proposed method. Such performance levels are far below the requirements of practical biometric systems. To provide a more meaningful comparison, the authors should also report open-set verification results under the intra-database setting, which is the standard evaluation protocol for vein recognition systems.}
\color{black}

\begin{itemize}
	\item \textbf{Authors' Response:} 
	We thank the reviewer for this important observation. We agree that cross-dataset evaluation can yield higher EER values due to the inherent domain shift between datasets, and that intra-database evaluation is a standard and complementary protocol for assessing practical biometric performance.

	First, we would like to clarify that the cross-dataset setting used in Table~5 is intentionally designed to evaluate {open-set generalisation under severe domain shift}, where both the acquisition conditions and subject identities differ between training and testing. In this scenario, higher EER values are expected for all methods, as the task is significantly more challenging than intra-database evaluation. Therefore, these results should be interpreted as a measure of robustness rather than absolute deployment performance.

	Following the reviewer’s suggestion, we have now included {intra-database open-set verification experiments} for three datasets (FV-300, FV-USM, and MMCBNU) in \textbf{Section 2 of supplementary material.} In this protocol, the first 80\% of sorted identities are used for model development and the remaining 20\% of identities are held out entirely for open-set evaluation. Within the development identities, the dataset-provided train/test image folders are used for optimisation and validation/model selection. During final evaluation, embeddings are extracted only for the held-out identities, and genuine/impostor scores are computed exhaustively from within-identity and cross-identity image pairs among those unseen subjects. The results show that the proposed method achieves lower EER and higher TAR across most operating points under intra-database conditions, achieving very strong performance on FV-300 and strong performance on FV-USM and MMCBNU.
    

	We note that intra-database experiments were not conducted on PolyU and VERA due to their limited number of samples per identity (4 and 2, respectively), which makes it infeasible to construct a reliable intra-database development and held-out open-set evaluation split for learning-based methods.

	\item \textbf{Authors' Action:} 
	We have added a new subsection in \textbf{Section~2 of supplementary material} reporting {intra-database open-set verification results}, including:
	\begin{itemize}
	    \item A new table presenting AUC, EER, and TAR@FAR with 95\% confidence intervals.
	    \item A clear description of the intra-database protocol (80\% development identities and 20\% fully held-out open-set test identities).
	    \item A discussion comparing intra-database and cross-dataset performance, highlighting the trade-off between accuracy and generalisation.
	\end{itemize}

	These additions provide a more complete and balanced evaluation, aligning with standard biometric evaluation practices while preserving the focus on cross-dataset generalisation.
\end{itemize}
%%%&&&&&&
%%------------------
\color{blue} 
%%%&&&&&&
\textbf{\textit{Comment no. \showReviewcounter:}} 
\textit{The ablation experiments in Section 4 (Tables 6–7 and Fig. 8) are limited in scope and depth. The study mainly evaluates different convolution kernel sizes, the number of GrapherBlocks, and three loss variants. However, the datasets used in the ablation experiments are not clearly described, making the results ambiguous. The proposed loss is not compared with strong baseline losses such as ArcFace or CosFace, and there is no analysis explaining why the proposed loss improves performance.}
\color{black}

\begin{itemize}
	\item \textbf{Authors' Response:} 
	We thank the reviewer for this constructive feedback. We agree that the initial presentation of the ablation study lacked clarity regarding the experimental protocol and could be strengthened in terms of comparisons and analysis.

	In the revised manuscript, we have addressed these concerns as follows:

	\textbf{(i) Clarification of ablation protocol:}  
	All ablation experiments are now explicitly stated to be conducted under the $ABDE \rightarrow C$ leave-one-dataset-out protocol, where FV-USM is used as the held-out dataset. This has been clearly specified at the beginning of Section~4 and consistently referenced across all ablation subsections to remove any ambiguity.

	\textbf{(ii) Stronger loss-function comparison:}  
	In addition to standard losses (MSE, NLL), we have extended the ablation study to include comparisons with representative strong angular-margin and biometric embedding losses, namely ArcFace, MagFace, and AdaFace. These losses provide a strong comparison set because ArcFace is a widely used angular-margin baseline, while MagFace and AdaFace are more recent biometric embedding objectives.

	\textbf{(iii) Improved analysis of the proposed loss:}  
	We have expanded the discussion to better explain why the proposed CAH loss improves performance. Specifically, we highlight that:
\begin{itemize}
    \item Standard losses (MSE/NLL) do not explicitly enforce discriminative embedding structure.
    \item Angular-margin losses primarily target angular decision margins and inter-class separation, whereas the proposed CAH loss explicitly combines prototype-based compactness with angular regularisation, leading to improved intra-class compactness and inter-class separability for open-set verification.
\end{itemize}

	\textbf{(iv) Sensitivity analysis:}  
	We have also included a sensitivity study of the balancing parameter $\beta$, demonstrating how the angular component contributes to performance and identifying the optimal operating point.

	Overall, the revised ablation study is now more comprehensive, clearly defined, and better aligned with the reviewer’s expectations.

	\item \textbf{Authors' Action:} 
	We have made the following changes in the revised manuscript:
	\begin{itemize}
	    \item Explicitly stated the ablation protocol ($ABDE \rightarrow C$) in Section~4 and ensured consistency across all subsections.
	    \item Added a new comparison table evaluating the proposed loss against ArcFace, MagFace, and AdaFace.
	    \item Expanded the discussion to provide a clearer interpretation of the performance gains from the proposed CAH loss.
	    \item Included a parameter sensitivity analysis for $\beta$ to justify the design choice.
	\end{itemize}
\end{itemize}
%%%&&&&&&
%%------------------


%%%%%%%%%%%%%%%%%%%%%%%%%%%%%%%%%%%%%%%%%%%%%%%%%%%%%%%%%%%%%%%%%%%%%%
\end{document}
